\section{Evaluation}

\subsection{Five-Phase Evaluation Plan}
\label{subsec:five_phase_plan}

The objective of the evaluation is to assess the performance of the FINDING NEMO neighbor discovery protocol and the
CHUALANH topology-repair mechanism under dynamic membership changes. However, the hybrid SWIMPHI failure detector was
not evaluated independently under controlled delay or packet loss due to the protocol not providing network fault
injection. Nevertheless, unit tests and observation during manual testing confirmed that the implementation correctly
invokes indirect helper nodes and detect node failure more accurately by reducing the number of false positives.

The evaluation primarily measures execution time for each phase, including topology formation, failure detection, node
joining and topology recovery. By running many automated tests, we can calculate the proportion of successes and
failures within the configured timeout. The evaluation proceeds through a predefined sequence of tests or phases,
called the Five-Phase Evaluation Plan. Table~\ref{tab:experimental_phases} introduces the main phases, what is being
executed during each phase, and what is being measured. The timeout of each phase is set to ensure that during cases
where it is impossible for the network to reach the end condition without external forces, such as nodes joining or
failing, then the test can terminate and move on to the recovery phase R or the next phase.
Table~\ref{tab:phase_transitions} shows the phase transitions based on the outcome.

\begin{table*}[htbp]
  \centering
  \caption{Summary of Experimental Phases}
  \label{tab:experimental_phases}
  \renewcommand{\arraystretch}{1.3}
  \begin{tabularx}{\textwidth}{
      C{1.5cm}
      Z
      Z
      Z
      C{1.8cm}
    }
    \toprule
    \textbf{Phase} & \textbf{Description}                                    & \textbf{Start}               & \textbf{End}                    & \textbf{Timeout} \\
    \midrule
    Phase 1        & Start initial nodes                                     & The final node is started    & The network reached convergence & 120s             \\
    Phase S1       & Network stabilization check                             & After phase 1                & Phase timeout                   & 10s              \\
    Phase 2.1      & One node is terminated                                  & After node termination       & One failure report is sent      & 120s             \\
    Phase 2.2      & One node is terminated                                  & After node termination       & All failure reports are sent    & 120s             \\
    Phase 3        & One node is terminated                                  & After node termination       & The network reached convergence & 180s             \\
    Phase S3       & Network stabilization check                             & After phase 3                & Phase timeout                   & 10s              \\
    Phase 4        & Multiple nodes join                                     & The final node is started    & The network reached convergence & 120s             \\
    Phase S4       & Network stabilization check                             & After phase 4                & Phase timeout                   & 10s              \\
    Phase 5        & Multiple nodes are terminated                           & The final node is terminated & The network reached convergence & 180s             \\
    Phase S5       & Network stabilization check                             & After phase 5                & Phase timeout                   & 10s              \\
    Phase R        & Optional, when a main phase failed to reach convergence & After previous phase timeout & The network reached convergence & 120s             \\
    Phase SR       & Network stabilization check                             & After phase R                & Phase timeout                   & 10s              \\
    \bottomrule
  \end{tabularx}
\end{table*}

\begin{table}[htbp]
  \centering
  \caption{Phase Transitions Based on Evaluation Outcomes}
  \label{tab:phase_transitions}
  \renewcommand{\arraystretch}{1.3}

  \begin{tabularx}{\columnwidth}{
      C{1.3cm}
      Z
      Z
      Z
    }
    \toprule
    \multirow{2}{*}{\textbf{Phase}} & \multirow{2}{*}{\textbf{Success}}                       & \multicolumn{2}{c}{\textbf{Non-success outcome}}                      \\
    \cmidrule(lr){3-4}
                                    &                                                         & \textbf{\(P \leq 4\)}                            & \textbf{\(P > 4\)} \\
    \midrule
    Phase 1                         & Phase S1                                                & Phase R1                                         & Terminates run     \\
    Phase S1                        & Phase 2.1/2.2/3                                         & Phase R1                                         & Terminates run     \\
    Phase R1                        & Phase SR                                                & Phase 2/3                                        & Terminates run     \\
    Phase 2                         & \multicolumn{3}{c}{ Continues until Phase 3 completes }                                                                         \\
    Phase 3                         & Phase S3                                                & Phase R3                                         & Terminates run     \\
    Phase S3                        & Phase 4                                                 & Phase R3                                         & Terminates run     \\
    Phase R3                        & Phase SR                                                & Phase 4                                          & Terminates run     \\
    Phase 4                         & Phase S4                                                & Phase R4                                         & Terminates run     \\
    Phase S4                        & Phase 5                                                 & Phase R4                                         & Terminates run     \\
    Phase R4                        & Phase SR                                                & Phase 5                                          & Terminates run     \\
    Phase 5                         & Phase S5                                                & Phase R5                                         & Terminates run     \\
    Phase S5                        & Ends run                                                & Phase R5                                         & Terminates run     \\
    Phase R5                        & Phase SR                                                & Ends run                                         & Terminates run     \\
    Phase SR                        & \multicolumn{3}{c}{ Continues to the next main phase }                                                                          \\
    \bottomrule
  \end{tabularx}

  \vspace{0.5em}
  \begin{minipage}{0.96\linewidth}
    \footnotesize
    \(P\) denotes the number of problematic nodes detected when a phase terminates. For a main or stabilization phase,
    Phase R is executed only when between one and four replaceable problematic nodes remain. If no replaceable node
    exists, or if more than four problematic nodes are detected, the run is terminated.

    Phase 2.1, Phase 2.2, and Phase 3 execute concurrently. Phase 2.1 and Phase 2.2 do not independently determine the next
    phase; the transition is determined by the outcome of Phase 3.

    If Phase R fails with at most four remaining problematic nodes, the evaluator continues to the next main phase. Phase
    SR does not trigger another recovery phase. Therefore, execution also continues when Phase SR fails.
  \end{minipage}
\end{table}

\subsubsection{Phase 1: Initial Topology Formation}

Phase 1 evaluates the initial topology formation. After the convergence, a stabilization phase observes the topology
for 10 seconds to verify that all nodes remain active and maintain the neighbor list appropriately with valid,
bi-directional connections.

\subsubsection{Phase 2 \& 3: The Termination of One Node}

Then, a node is randomly chosen to be terminated. The next three phases runs in parallel. Phase 2.1 measures the time
from the termination to the first failure report is sent to the dashboard. Phase 2.2 measures until all neighbors of
the terminated node have sent at least one failure report each. Phase 3 measures until the network has returned to a
convergence state after the termination.

\subsubsection{Phase 4: Multiple Nodes Joining}

Phase 4 introduces additional nodes and measures the time required for the expanded topology to adapt and restore the
desired neighbor relationships. This phase aims to measure the effectiveness of FINDING NEMO on an converged,
established network.

\subsubsection{Phase 5: Multiple Nodes Terminating}

Phase 5 then evaluates the framework under multiple node failures by measuring the time required to detect all injected
failures, remove stale neighbor references, and restore the required neighbor count among all surviving nodes.

\subsubsection{Phase S: Stabilization Check}

Phase S - stabilization - is executed after each time the network topology successfully converges. This includes phases
1, 3, 4, 5 and R. The main purpose is to make sure that the network can keep the convergence state for a certain period
of time, in this case, 10 seconds, during which the phase continuously monitoring the network to check for any
anomalies and report them.

\subsubsection{Phase R: Recovery Phase}

Phase R is an optional recovery phase that is triggered only when a preceding phase fails to satisfy its completion
criteria within the configured timeout period while the network containing 4 or less problematic nodes. Problematic
nodes includes deficient nodes (nodes that have less than the required maximum number of neighbors), nodes with stale
edges which reference inactive, killed or unknown nodes, or nodes with a non-mutual neighbor relationship. During phase
R, all these nodes are terminated, and the same amount of new nodes is created. This is to simulate the node self
restarting when they are unable to rewire after a certain amount of time. The reasoning for this will be discussed in
Section~\ref{sec:conclusion}.

\subsection{Experimental Environment and Configuration}
\label{subsec:experimental_environment}

The majority of the automated evaluation runs were conducted on a Lenovo 20TJS1HN0V laptop running Microsoft Windows~11
Pro, build~26200. The machine was equipped with an Intel Core i7-10850H processor operating at a base frequency of
approximately \(2.70\,\mathrm{GHz}\), with six physical cores and twelve logical processors. It contained
\(32\,\mathrm{GB}\) of DDR4-3200 main memory and two \(512\,\mathrm{GB}\) NVMe solid-state drives, providing
approximately \(1\,\mathrm{TB}\) of total marketed storage capacity. The active wireless interface was an Intel Wi-Fi~6
AX201 adapter with a reported link speed of \(866.7\,\mathrm{Mbps}\).

The software environment consisted of Oracle Java~25.0.1 LTS, and Apache Maven~3.9.11. During each evaluation run, the
dashboard server, evaluation controller, and all distributed node agents were executed concurrently. Each node agent
ran as an independent Java process and exposed a separate peer-to-peer communication port. Communication between node
processes was performed through HTTP using the local network interface.

The dashboard server was responsible only for collecting heartbeat reports, storing failure information, visualizing
the network state, and providing observations to the evaluation controller. The decentralized failure-detection and
topology-repair mechanisms operated within the node agents and did not depend on the dashboard for their coordination.

Although the majority of runs were performed using the primary machine described above, additional runs were conducted
on other laptops to increase the number of recorded experiments. These machines used the same source code, Java runtime
requirements, phase definitions, timeout values, and evaluation procedure, but differed in their hardware
configurations. Because the resulting observations were aggregated and were not grouped according to the machine on
which they were executed, the stated hardware configuration should be interpreted as representative of the majority of
the evaluation rather than as the exclusive experimental environment.

The experiments evaluated multiple combinations of initial node count, maximum neighbor limit, number of newly joining
nodes, and number of simultaneously terminated nodes. The configured maximum neighbor limit, denoted by \(k\), was
shared by all nodes within one run and was restricted to even values because FINDING NEMO and CHUALANH rely on pairing
and rewiring operations that preserve the target node degree.

\subsection{Evaluation Results}
\label{subsec:evaluation_results}

\subsubsection{Definitions}
\label{subsubsec:evaluation_definitions}

The evaluation uses the following definitions to determine the state and outcome of each experimental phase.

\textbf{Stale edge.}
A stale edge exists when an active node reports a neighbor that is no
longer part of the expected active-node set. This includes references to
terminated, inactive, or unknown nodes that should already have been
removed from the local membership view.

\textbf{Non-mutual edge.}
A neighbor relationship between nodes \(u\) and \(v\) is non-mutual when
one node reports the other as a neighbor but the reverse relationship is
absent. In other words,
\[
  v \in N_u
  \quad\text{but}\quad
  u \notin N_v,
\]
or vice versa. Since neighbor relationships in the proposed overlay are required to be bidirectional, a non-mutual edge
represents an invalid topology state.

\textbf{Problematic node.}
A problematic node is an active node that is involved in at
least one topology violation. A node is therefore classified as
problematic when it satisfies one or more of the following conditions:

\begin{itemize}
  \item its number of distinct neighbors is not equal to \(k\);
  \item it contains duplicate entries in its reported neighbor list;
  \item it contains itself as a neighbor;
  \item it contains a stale edge to a terminated, inactive, or unknown node; or
  \item it is involved in a non-mutual neighbor relationship.
\end{itemize}

A node that violates several conditions is counted only once.

\textbf{Network convergence.}
The network is considered converged when every active node is
correctly reported as active and non-problematic.
Network convergence therefore requires both the correct neighbor count and a structurally valid topology. Having \(k\)
reported neighbor entries alone is not sufficient when those entries include stale, duplicate, self-referential, or
non-mutual edges.

\textbf{Timeout.}
A timeout occurs when a phase does not satisfy its completion condition
within its configured time limit. A timeout does not necessarily imply
that convergence is permanently impossible; it indicates that the
required state was not reached within the time considered acceptable by
the evaluation procedure.

\subsubsection{Overview}

Table~\ref{tab:valid_runs_by_phase} shows the number of recorded runs following the Five-Phase Evaluation Plan. Due to
multiple reasons such as the user terminating the evaluation prematurely, or a phase failure which can lead to
termination of the current run, the number of executions of a later phase can be less than the previous phases.
Overall, more than 200 runs were executed

\begin{table}[H]
  \centering
  \caption{Number of Valid Runs Recorded}
  \label{tab:valid_runs_by_phase}
  \begin{tabular}{cc}
    \toprule
    \textbf{Phase} & \textbf{Number of Runs} \\
    \midrule
    Phase 1        & 236                     \\
    Phase 2.1      & 235                     \\
    Phase 2.2      & 234                     \\
    Phase 3        & 233                     \\
    Phase 4        & 228                     \\
    Phase 5        & 221                     \\
    \bottomrule
  \end{tabular}
\end{table}

Figure~\ref{fig:final_success_rate_by_phase} illustrates the percentages of successes and failures categorized by each
phase. The success rates fluctuate across phases 1 through 4, and they remain consistent at around 95\%. This shows
that the protocol performs reliably in network formation, node addition and single node failure detection and recovery.

In constrast, phase 5 recorded a noticably lower success rate of only 55.9\%, with 37.7\% of the runs failed. This
suggests that the implementation of the protocol failed to detect and recovery when multiple nodes failed at the same
time. When several node terminate at the same time, many nodes become deficient simultaneously and a burst of recovery
messages is sent out during a sort period of time. The increased communication load and the presence of multiple
concurrent recovery sessions may lead to conflicting operations, delayed convergence, or recovery timeouts.

\begin{figure}[h]
  \centering
  \includegraphics[width=0.9\linewidth]{final-success-rate-chart-1.pdf}
  \caption{Success Rate by Main Evaluation Phase}
  \label{fig:final_success_rate_by_phase}
\end{figure}

A closer examination of the detailed outcomes in Figure~\ref{fig:detailed-outcome-chart} reveals that, in phases~1
and~4, a substantial proportion of the successful runs required the additional recovery Phase~R. When one or more nodes
initially join the system, they use the FINDING NEMO algorithm to discover and establish connections with neighboring
nodes. Each joining node independently creates and executes its own join plan. Consequently, multiple join plans may
conflict with one another, preventing some plans from being fully executed because the involved nodes may reject
conflicting requests.

Such conflicts do not necessarily result in permanent failure, as deficient nodes subsequently initiate the recovery
process using the CHUALANH algorithm. In many cases, this recovery phase allows the network to reach a valid and stable
topology. However, the CHUALANH algorithm also has structural limitations that may create situations in which network
convergence is impossible. These limitations and their implications are discussed further in
Section~\ref{sec:conclusion}.

\begin{figure}[h]
  \centering
  \includegraphics[width=0.9\linewidth]{detailed-outcome-chart-1.pdf}
  \caption{Detailed Success Rate by Main Evaluation Phase}
  \label{fig:detailed-outcome-chart}
  \vspace{0.3em}
  \begin{minipage}{0.96\linewidth}
    \footnotesize
    \raggedright
    \textit{Success} means that the phase succeeds directly without requiring Phase R.

    \textit{Success after recovery} means that the phase initially times out, but Phase R successfully recovers the network.

    \textit{Failure without recovery} means that the phase times out without triggering Phase R.

    \textit{Recovery failed} means that Phase R is executed but cannot recover the network.
  \end{minipage}
\end{figure}

The high proportion of runs recovered successfully through Phase~R is further supported by
Figure~\ref{fig:recovery-success-chart}.

\begin{figure}[h]
  \centering
  \includegraphics[width=0.9\linewidth]{recovery-success-chart-1.pdf}
  \caption{Success Rate of Recovery Attempts by Triggering Phase}
  \label{fig:recovery-success-chart}
\end{figure}

\subsubsection{Phase 1}
\label{sec:phase-1-evaluation}

Figure~\ref{fig:phase1-results} presents the results of Phase~1 from three perspectives: overall success rates,
execution time and the number of problematic nodes in failed runs.

The success rate chart in Figure~\ref{fig:phase1-input-chart} shows the percentage of successful runs based on the
initial network size and the neighbor limit, and the size of each bubble represents the total number of runs with that
specific configuration. Overall, the success rates are quite high, with most of them reach 90\% or higher, and only one
case with 86\% when there are 14 initial nodes and the number of maximum neighbors required is 4. The results showed no
clear relationship between the initial conditions and the success rates, rather the lower success rate is likely caused
by the machine on which these tests were conducted.

Figure~\ref{fig:phase1-time-impact-chart} illustrates the execution time of each run, not including the recovery phase
R. As previously shown in Figure~\ref{fig:detailed-outcome-chart}, almost half of the runs timed out and required the
recovery phase to reach convergence. But for the successful runs, the convergence time is far below the configured
timeout, averaging at about 18 seconds. Higher initial node counts tends to require a slightly more time to converges,
sometimes up to 60 seconds. Nevertheless, this trend appears relatively weak and does not significantly affect the
overall performance of the protocol.

Further insight is provided by Figure~\ref{fig:phase1-problem-count-chart}, which summarizes the number of problematic
nodes remaining in failed runs. Most failures involved 4 nodes or less, with most of them have only 1 or 2 problematic
nodes. Having a few problematic nodes after several nodes joining is actually not a problem for the system, and can be
detected and fixed easily. This will be discussed in detail in Section~\ref{sec:conclusion}.

From all of these illustrations, we concluded that the network performed extremely well during network establishment. A
small percentage of runs failed, but further testing have confirmed that if the network is allowed to keep running,
overtime as nodes join and exit the network, the state of convergence will be achieved.

\begin{figure}[t]
  \centering

  \subfloat[Success rate]{
    \includegraphics[width=0.9\linewidth]{phase1-input-chart-1.pdf}
    \label{fig:phase1-input-chart}
  }
  \hfill
  \subfloat[Execution time]{
    \includegraphics[width=0.9\linewidth]{phase1-time-impact-chart-1.pdf}
    \label{fig:phase1-time-impact-chart}
  }
  \hfill
  \subfloat[Problematic nodes]{
    \includegraphics[width=0.9\linewidth]{phase1-problem-count-chart-1.pdf}
    \label{fig:phase1-problem-count-chart}
  }

  \caption{Phase 1 evaluation results}
  \label{fig:phase1-results}
\end{figure}

\subsubsection{Phase 2}
\label{sec:phase-2-evaluation}

In phase 2, one node is killed and the times taken to receive one report and all expected reports from the dead node's
neighbors are measured.

We can observe from Figure~\ref{fig:phase2-time-chart} that the first report is received quite fast, only take just
more than 10 seconds. The time taken to receive all reports is quite uniformly distributed from 15 to 50 seconds. There
is a small number of timeouts, however the recorded logs revealed the reason for these timeouts. For all of these
timeouts, phase~3 has been completed much earlier. This means that all nodes has registered locally the information
about the killed node, either from pinging the node itself or from a gossip. If a gossip is received, the node will not
send a failure report to the dashboard. Therefore, phase 2.2 condition is not met and thus timeout.

These results showed that with a single node failure, the system correctly detects it 100\% of the time in all the
conducted tests. This proves the reliability of the failure detection system, however does not prove that the SWIMPHI
protocol can prevent false positives as mentioned in part~\ref{subsec:five_phase_plan}.

\begin{figure}
  \centering
  \includegraphics[width=0.9\linewidth]{phase2-time-chart-1.pdf}
  \caption{Phase 2: Detection Time for First Report and All Expected Reporting}
  \label{fig:phase2-time-chart}
\end{figure}

\subsubsection{Phase 3}
\label{sec:phase-3-evaluation}

Phase 3 measures the time until the network has recovered from a single node failure. The results in
Figure~\ref{fig:phase3-results} are presented in the same way as Phase~1 in Section~\ref{sec:phase-1-evaluation}.

In Figure~\ref{fig:phase3-input-chart}, we can observe that the success rates are extremely high, almost all cases were
successful. Only a few tests failed to recover within the allowed time.

Interestingly, Figure~\ref{fig:phase3-time-impact-chart} shows that as the maximum number of neighbor increases, more
time is needed to recover from the node deletion. However, most timeouts are caused during tests with low initial node
counts and low neighbor limit.

\begin{figure}[t]
  \centering

  \subfloat[Success rate]{
    \includegraphics[width=0.9\linewidth]{phase3-input-chart-1.pdf}
    \label{fig:phase3-input-chart}
  }
  \hfill
  \subfloat[Execution time]{
    \includegraphics[width=0.9\linewidth]{phase3-time-impact-chart-1.pdf}
    \label{fig:phase3-time-impact-chart}
  }
  \hfill
  \subfloat[Problematic nodes]{
    \includegraphics[width=0.9\linewidth]{phase3-problem-count-chart-1.pdf}
    \label{fig:phase3-problem-count-chart}
  }

  \caption{Phase 3 evaluation results}
  \label{fig:phase3-results}
\end{figure}

\subsubsection{Phase 4}
\label{sec:phase-4-evaluation}

Similar to phase 1, when multiple nodes join the network, the success rate is quite high as illustrated in
Figure~\ref{fig:phase4-input-chart}, but almost half of the runs required the execution of phase R
(Figure~\ref{fig:detailed-outcome-chart}). Figure~\ref{fig:phase4-time-impact-chart} further proves this point, as most
of the runs timed out. The average run converges after around 17 seconds on average. Most runs start the recovery phase
with 1 or 2 problematic nodes, which is ideal for the system and generates high success rates as seen in
Figure~\ref{fig:recovery-success-chart}.

\begin{figure}[t]
  \centering

  \subfloat[Success rate]{
    \includegraphics[width=0.9\linewidth]{phase4-input-chart-1.pdf}
    \label{fig:phase4-input-chart}
  }
  \hfill
  \subfloat[Execution time]{
    \includegraphics[width=0.9\linewidth]{phase4-time-impact-chart-1.pdf}
    \label{fig:phase4-time-impact-chart}
  }
  \hfill
  \subfloat[Problematic nodes]{
    \includegraphics[width=0.9\linewidth]{phase4-problem-count-chart-1.pdf}
    \label{fig:phase4-problem-count-chart}
  }

  \caption{Phase 4 evaluation results}
  \label{fig:phase4-results}
\end{figure}

\subsubsection{Phase 5}
\label{sec:phase-5-evaluation}

Phase 5 proves to be the most difficult phase for the system.

From Figure~\ref{fig:phase5-input-chart}, we can see that for the three configurations: remove 4, 5, 6 nodes when the
maximum neighbors is 6, 4 and 4 nodes respectively, the success rate is very high. However, for the rest of the
configurations, the tests were not very promising. It seems that the higher the number of failures and the maximum
neighbors are, the lower the success rates are. This is consistent with the protocol design. Each failure creates
deficient nodes. The more failures, the more deficient nodes. When too many nodes fail at the same time, most of the
nodes are deficient. This causes all nodes to send out recovery messages during a short period of time. Then all
recovery sessions overlapping, and competition for available neighbors increase the delays and prevent recovery
algorithms from completing successfully.

The execution times shown in Figure~\ref{fig:phase5-time-impact-chart} provide further insight into this behavior.
Successful runs generally completed before the configured timeout, whereas unsuccessful runs frequently reached the
timeout threshold. This suggests that the main problem was not merely slower recovery, but that some configurations
were unable to reach convergence within the evaluation period.

Figure~\ref{fig:phase5-problem-count-chart} is further proof for this. A lot of runs ended up with more than 4
problematic nodes, which according to the five-phase evaluation plan, will not trigger the recovery phase and end the
test immediately. % remember to talk about this in the plan description.

To sum up, the current implementation of the protocol does not work well when multiple nodes fail at the same time.
This can be due to the massive amount of messages being sent during the recovery phase, and thus a protocol or
implementation issue. This can also be a product of running a massive amount of Java processes on a single laptop. We
can not conclude without enough resources for further, large-scale testing.

\begin{figure}[t]
  \centering

  \subfloat[Success rate]{
    \includegraphics[width=0.9\linewidth]{phase5-input-chart-1.pdf}
    \label{fig:phase5-input-chart}
  }
  \hfill
  \subfloat[Execution time]{
    \includegraphics[width=0.9\linewidth]{phase5-time-impact-chart-1.pdf}
    \label{fig:phase5-time-impact-chart}
  }
  \hfill
  \subfloat[Problematic nodes]{
    \includegraphics[width=0.9\linewidth]{phase5-problem-count-chart-1.pdf}
    \label{fig:phase5-problem-count-chart}
  }

  \caption{Phase 5 evaluation results}
  \label{fig:phase5-results}
\end{figure}
